%! Parametrization of Implicitly Defined Functions — Unit 4, Lesson 4.7 — Warm-Up (Answer Key)
\documentclass[10pt]{article}
\usepackage{precalculus-article}
\usepackage{precalculus-key}   % pulls in -boxes; do NOT also load -boxes
\newcommand{\TallMath}[1]{$\displaystyle #1\rule[-1.4em]{0pt}{3.2em}$}

\begin{document}

\pageheader{Unit 4, Lesson 4.7}{Warm-Up --- Answer Key}
\namedateperiod

\vspace{0.12in}

\begin{notesbox}{Spiral Review --- Key}
\small

\textbf{1.} Recall the parametrization of the unit circle $x^2 + y^2 = 1$.
Fill in the component functions and domain.

\vspace{0.08in}
\hspace{1em} $x(t) = $ \ans{$\cos t$} \qquad $y(t) = $ \ans{$\sin t$} \qquad
domain: \ans{$0 \le t \le 2\pi$}

\vspace{0.18in}

\textbf{2.} For the ellipse $\dfrac{x^2}{25} + \dfrac{y^2}{9} = 1$,
find the $x$-intercepts and $y$-intercepts.

\vspace{0.08in}
\hspace{1em} $x$-intercepts: \ans{$(\pm 5,\, 0)$} \qquad $y$-intercepts: \ans{$(0,\, \pm 3)$}

\vspace{0.18in}

\textbf{3.} Let $x(t) = 5\cos t$ and $y(t) = 3\sin t$.
Complete the table, then verify that one of the points satisfies
$\dfrac{x^2}{25} + \dfrac{y^2}{9} = 1$.

\vspace{0.1in}
\renewcommand{\arraystretch}{1.6}
\hspace{1em}
\begin{tabular}{|c|c|c|l|}
\hline
$t$ & $x(t)$ & $y(t)$ & Point $(x,y)$ \\
\hline
$0$ & \ans{$5$} & \ans{$0$} & \ans{$(5,\,0)$} \\
\hline
$\pi/2$ & \ans{$0$} & \ans{$3$} & \ans{$(0,\,3)$} \\
\hline
$\pi$ & \ans{$-5$} & \ans{$0$} & \ans{$(-5,\,0)$} \\
\hline
$3\pi/2$ & \ans{$0$} & \ans{$-3$} & \ans{$(0,\,-3)$} \\
\hline
\end{tabular}

\vspace{0.1in}
\hspace{1em} Verify (using $t=\pi/2$, point $(0,3)$): \ans{$\dfrac{0^2}{25}+\dfrac{3^2}{9}=0+1=1$. \checkmark}

\end{notesbox}

\begin{teachernote}
All four table points should be recognized as the axis intercepts of the ellipse --- this bridges
directly to the lesson's main idea. If students are confused by Q3, ask: ``Did you recognize these
points from Q2?'' The verification step builds fluency with substituting parametric outputs into
implicit equations, which is the core skill of EK 4.7.A.2.
\end{teachernote}

\end{document}
